\documentclass{article}
\usepackage{amsmath,amssymb,amsthm}
\usepackage{algorithm}
\usepackage{algpseudocode}
\usepackage{bm}
\usepackage{hyperref}
\usepackage{enumitem}
\newtheorem{theorem}{Theorem}
\newtheorem{lemma}{Lemma}
\newtheorem{assumption}{Assumption}
\newcommand{\EE}{\mathbb{E}}
\newcommand{\RR}{\mathbb{R}}
\newcommand{\norm}[1]{\left\|#1\right\|}
\newcommand{\inner}[1]{\left\langle #1 \right\rangle}
\begin{document}

\section*{Stochastic Subsampled Newton (SSN) Algorithm}

\section*{Mathematical Formulation}
Let $f:\RR^{d}\to\RR$ be a twice differentiable (possibly non‑convex) objective written as a finite sum
\[
f(x)=\frac{1}{n}\sum_{i=1}^{n} f_i(x),\qquad 
f_i:\RR^{d}\to\RR .
\]
At iteration $t$ we draw a random mini‑batch $\mathcal{B}_t\subset\{1,\dots,n\}$ of size $b$ and form the stochastic gradient
\[
g_t\;:=\;\frac{1}{b}\sum_{i\in\mathcal{B}_t}\nabla f_i(x_t),
\]
and the stochastic Hessian estimate
\[
H_t\;:=\;\frac{1}{b}\sum_{i\in\mathcal{B}_t}\nabla^{2} f_i(x_t) + \lambda I_d ,
\]
where $\lambda>0$ is a regularization parameter guaranteeing invertibility.  
The update rule is
\[
\boxed{%
x_{t+1}=x_t-\alpha\, H_t^{-1} g_t
}\tag{1}
\]
with a stepsize (learning rate) $\alpha\in(0,1]$.

\section*{Pseudocode}
\begin{algorithm}[H]
\caption{Stochastic Subsampled Newton (SSN)}
\begin{algorithmic}[1]
\State \textbf{Input:} initial point $x_0\in\RR^{d}$, stepsize $\alpha\in(0,1]$, regularizer $\lambda>0$, batch size $b$, max iterations $T$.
\For{$t=0,1,\dots,T-1$}
    \State Sample a mini‑batch $\mathcal{B}_t\subset\{1,\dots,n\}$, $|\mathcal{B}_t|=b$.
    \State Compute stochastic gradient 
    \[
    g_t\gets \frac{1}{b}\sum_{i\in\mathcal{B}_t}\nabla f_i(x_t).
    \]
    \State Compute stochastic Hessian estimate 
    \[
    \widetilde H_t\gets \frac{1}{b}\sum_{i\in\mathcal{B}_t}\nabla^{2} f_i(x_t).
    \]
    \State Form regularized matrix $H_t\gets \widetilde H_t+\lambda I_d$.
    \State Update $x_{t+1}\gets x_t-\alpha H_t^{-1} g_t$.
\EndFor
\State \textbf{Output:} $x_T$
\end{algorithmic}
\end{algorithm}

\section*{Dry Run Example}
Consider the one‑dimensional quadratic $f(w)=(w-3)^2$, thus $\nabla f(w)=2(w-3)$ and $\nabla^{2}f(w)=2$ (deterministic, so stochastic estimates coincide with the true values).  
Set $w_0=0$, stepsize $\alpha=0.1$, regularizer $\lambda=10^{-3}$ (practically negligible), and use a batch size $b=1$ (the whole dataset).

\begin{tabular}{c|c|c|c}
Iteration $t$ & $g_t= \nabla f(w_t)$ & $H_t= \nabla^{2}f(w_t)+\lambda$ & $w_{t+1}=w_t-\alpha H_t^{-1}g_t$\\ \hline
0 & $2(0-3)=-6$ & $2+10^{-3}\approx2.001$ & $0-0.1\cdot\frac{-6}{2.001}\approx0.2999$\\
1 & $2(0.2999-3)\approx-5.4002$ & $2.001$ & $0.2999-0.1\cdot\frac{-5.4002}{2.001}\approx0.5699$\\
2 & $2(0.5699-3)\approx-4.8602$ & $2.001$ & $0.5699-0.1\cdot\frac{-4.8602}{2.001}\approx0.8129$
\end{tabular}

The objective values are $f(0)=9$, $f(0.2999)\approx7.29$, $f(0.5699)\approx5.93$, $f(0.8129)\approx4.80$, showing descent toward the optimum $w^\star=3$.

\section*{Assumptions}
\begin{assumption}[Smoothness]\label{ass:smooth}
$f$ has $L$‑Lipschitz gradients:
\[
\forall x,y\in\RR^{d},\qquad 
\norm{\nabla f(x)-\nabla f(y)}\le L\norm{x-y}.
\]
Consequently $\forall x$, $\nabla^{2}f(x)\preceq L I_d$.
\end{assumption}
\begin{assumption}[Bounded Hessian Spectrum]\label{ass:hessian}
There exist constants $0<\mu\le L$ such that for every $x$,
\[
\mu I_d \preceq \nabla^{2} f(x) \preceq L I_d .
\]
\end{assumption}
\begin{assumption}[Unbiased Stochastic Gradient]\label{ass:unbiased-grad}
\[
\EE[g_t\mid x_t]=\nabla f(x_t),\qquad 
\EE\!\left[\norm{g_t-\nabla f(x_t)}^{2}\mid x_t\right]\le\sigma_g^{2}.
\]
\end{assumption}
\begin{assumption}[Unbiased Stochastic Hessian]\label{ass:unbiased-hess}
\[
\EE[\widetilde H_t\mid x_t]=\nabla^{2} f(x_t),\qquad 
\EE\!\left[\norm{\widetilde H_t-\nabla^{2} f(x_t)}_{F}^{2}\mid x_t\right]\le\sigma_H^{2}.
\]
\end{assumption}
\begin{assumption}[Regularization]\label{ass:regular}
The regularization $\lambda>0$ satisfies $\lambda\le \mu$ so that
\[
\lambda_{\min}(H_t)\ge \mu+\lambda\ge \mu>0 .
\]
\end{assumption}

\section*{Main Convergence Theorem}
\begin{theorem}[Expected Gradient Norm Decrease]\label{thm:main}
Under Assumptions \ref{ass:smooth}–\ref{ass:regular} and for a constant stepsize
\[
\alpha\le \frac{2\mu}{L^{2}+ \mu^{2}},
\]
the iterates generated by \eqref{1} satisfy
\[
\frac{1}{T}\sum_{t=0}^{T-1}\EE\!\left[\norm{\nabla f(x_t)}^{2}\right]
\;\le\;
\frac{2\bigl(f(x_0)-f^\star\bigr)}{\alpha\,\mu\,T}
\;+\;
\frac{\alpha\,\sigma_g^{2}}{\mu}
\;+\;
\frac{L\,\sigma_H^{2}}{\mu^{2}} .
\tag{2}
\]
Consequently, choosing $\alpha=O(T^{-1/2})$ yields the rate
\[
\frac{1}{T}\sum_{t=0}^{T-1}\EE\!\left[\norm{\nabla f(x_t)}^{2}\right]=O\!\left(T^{-1/2}\right).
\]
\end{theorem}

\section*{Theoretical Properties}
\begin{itemize}[leftmargin=*]
\item \textbf{Stability.} The regularizer $\lambda$ guarantees $H_t\succ0$ uniformly, preventing blow‑up of $H_t^{-1}$.
\item \textbf{Hyper‑parameter Sensitivity.} The bound on $\alpha$ depends on the curvature ratio $\kappa=L/\mu$. Larger $\kappa$ forces smaller stepsizes, reflecting the usual Newton‑type sensitivity.
\item \textbf{Comparison to SGD.} When $b=n$ (full‑batch) the variance terms $\sigma_g^{2},\sigma_H^{2}$ vanish and the bound reduces to the deterministic Newton descent rate $O(1/T)$. For $b\ll n$, the additional variance terms appear, matching the $O(T^{-1/2})$ rate typical of stochastic first‑order methods, but with a potentially much smaller constant because curvature information accelerates progress.
\end{itemize}

\section*{Discussion}
The theorem shows that even with noisy Hessian information the SSN method drives the expected squared gradient norm to zero at a sub‑linear rate. The bound separates the contribution of gradient noise ($\sigma_g^{2}$) and Hessian noise ($\sigma_H^{2}$). In the regime where the Hessian variance is negligible (large batch for curvature), the method enjoys a faster $O(1/T)$ decay, interpolating between pure Newton and pure SGD behaviour.

\section*{Preliminaries (Definitions and Lemmas)}
\begin{lemma}[Quadratic Upper Bound]\label{lem:quad}
For an $L$‑smooth function,
\[
f(y)\le f(x)+\inner{\nabla f(x)}{y-x}+\frac{L}{2}\norm{y-x}^{2},\qquad\forall x,y.
\]
\end{lemma}
\begin{lemma}[Matrix Inverse Bound]\label{lem:inv}
If $A\succeq \mu I_d$ then $\norm{A^{-1}}\le 1/\mu$ and for any $u$,
\[
\inner{A^{-1}u}{u}\ge \frac{1}{L}\norm{u}^{2}.
\]
\end{lemma}
Both lemmas are standard; proofs are omitted for brevity.

\section*{Main Proof (Step‑by‑Step Derivation)}
We prove Theorem \ref{thm:main}. All expectations are taken w.r.t.\ the randomness at iteration $t$ conditioned on $x_t$ unless otherwise noted.

\noindent\textbf{[Step 1]} Apply Lemma \ref{lem:quad} with $y=x_{t+1}=x_t-\alpha H_t^{-1}g_t$:
\[
f(x_{t+1})\le f(x_t)-\alpha\inner{ \nabla f(x_t)}{H_t^{-1}g_t}
+\frac{L\alpha^{2}}{2}\norm{H_t^{-1}g_t}^{2}. \tag{3}
\]

\noindent\textbf{[Step 2]} Decompose $g_t$ into its expectation plus noise:
\[
g_t=\nabla f(x_t)+\xi_t,\qquad \EE[\xi_t\mid x_t]=0,\;
\EE[\norm{\xi_t}^{2}\mid x_t]\le\sigma_g^{2}. \tag{4}
\]

\noindent\textbf{[Step 3]} Substitute \eqref{4} into the inner product term of \eqref{3}:
\[
\inner{\nabla f(x_t)}{H_t^{-1}g_t}
= \inner{\nabla f(x_t)}{H_t^{-1}\nabla f(x_t)}
+ \inner{\nabla f(x_t)}{H_t^{-1}\xi_t}. \tag{5}
\]

\noindent\textbf{[Step 4]} Take conditional expectation of the second term in \eqref{5}. By unbiasedness of $\xi_t$,
\[
\EE\!\left[\inner{\nabla f(x_t)}{H_t^{-1}\xi_t}\mid x_t\right]=
\inner{\nabla f(x_t)}{H_t^{-1}\EE[\xi_t\mid x_t]}=0. \tag{6}
\]

\noindent\textbf{[Step 5]} For the quadratic term in \eqref{3}, expand the norm:
\[
\norm{H_t^{-1}g_t}^{2}
= \norm{H_t^{-1}\nabla f(x_t)}^{2}
+2\inner{H_t^{-1}\nabla f(x_t)}{H_t^{-1}\xi_t}
+\norm{H_t^{-1}\xi_t}^{2}. \tag{7}
\]

\noumber\textbf{[Step 6]} Conditioned on $x_t$, the cross term vanishes:
\[
\EE\!\left[\inner{H_t^{-1}\nabla f(x_t)}{H_t^{-1}\xi_t}\mid x_t\right]=
\inner{H_t^{-1}\nabla f(x_t)}{H_t^{-1}\EE[\xi_t\mid x_t]}=0. \tag{8}
\]

\noindent\textbf{[Step 7]} Take conditional expectation of \eqref{7} using \eqref{8} and the variance bound:
\[
\EE\!\left[\norm{H_t^{-1}g_t}^{2}\mid x_t\right]
= \norm{H_t^{-1}\nabla f(x_t)}^{2}
+\EE\!\left[\norm{H_t^{-1}\xi_t}^{2}\mid x_t\right]
\le \norm{H_t^{-1}\nabla f(x_t)}^{2}
+\frac{\sigma_g^{2}}{\mu^{2}}, \tag{9}
\]
where we used Lemma \ref{lem:inv} to bound $\norm{H_t^{-1}}\le 1/\mu$.

\noindent\textbf{[Step 8]} Plug \eqref{5}, \eqref{6}, and \eqref{9} into \eqref{3} and then take conditional expectation:
\[
\EE\!\bigl[f(x_{t+1})\mid x_t\bigr]
\le f(x_t)-\alpha\inner{\nabla f(x_t)}{H_t^{-1}\nabla f(x_t)}
+\frac{L\alpha^{2}}{2}\norm{H_t^{-1}\nabla f(x_t)}^{2}
+\frac{L\alpha^{2}\sigma_g^{2}}{2\mu^{2}}. \tag{10}
\]

\noindent\textbf{[Step 9]} Use Lemma \ref{lem:inv} to bound the quadratic forms:
\[
\inner{\nabla f(x_t)}{H_t^{-1}\nabla f(x_t)}\ge \frac{1}{L}\norm{\nabla f(x_t)}^{2},
\qquad
\norm{H_t^{-1}\nabla f(x_t)}^{2}\le \frac{1}{\mu^{2}}\norm{\nabla f(x_t)}^{2}. \tag{11}
\]

\noindent\textbf{[Step 10]} Substitute \eqref{11} into \eqref{10}:
\[
\EE\!\bigl[f(x_{t+1})\mid x_t\bigr]
\le f(x_t)-\alpha\frac{1}{L}\norm{\nabla f(x_t)}^{2}
+\frac{L\alpha^{2}}{2}\frac{1}{\mu^{2}}\norm{\nabla f(x_t)}^{2}
+\frac{L\alpha^{2}\sigma_g^{2}}{2\mu^{2}}. \tag{12}
\]

\noindent\textbf{[Step 11]} Rearrange the deterministic terms:
\[
\EE\!\bigl[f(x_{t+1})\mid x_t\bigr]
\le f(x_t)-\Bigl(\frac{\alpha}{L}-\frac{L\alpha^{2}}{2\mu^{2}}\Bigr)\norm{\nabla f(x_t)}^{2}
+\frac{L\alpha^{2}\sigma_g^{2}}{2\mu^{2}}. \tag{13}
\]

\noindent\textbf{[Step 12]} Choose $\alpha$ such that the coefficient of $\norm{\nabla f(x_t)}^{2}$ is non‑negative. The condition
\[
\frac{\alpha}{L}-\frac{L\alpha^{2}}{2\mu^{2}}\ge 0
\quad\Longleftrightarrow\quad
\alpha\le \frac{2\mu^{2}}{L^{2}}=\frac{2\mu}{L^{2}/\mu}
\]
is satisfied by the stepsize bound stated in the theorem. Under this choice we obtain
\[
\EE\!\bigl[f(x_{t+1})\mid x_t\bigr]
\le f(x_t)-\frac{\alpha\mu}{2L^{2}}\norm{\nabla f(x_t)}^{2}
+\frac{L\alpha^{2}\sigma_g^{2}}{2\mu^{2}}. \tag{14}
\]

\noindent\textbf{[Step 13]} Take full expectation and sum over $t=0,\dots,T-1$:
\[
\EE\!\bigl[f(x_T)\bigr]
\le f(x_0)-\frac{\alpha\mu}{2L^{2}}\sum_{t=0}^{T-1}\EE\!\left[\norm{\nabla f(x_t)}^{2}\right]
+\frac{L\alpha^{2}\sigma_g^{2}}{2\mu^{2}}T. \tag{15}
\]

\noindent\textbf{[Step 14]} Since $f(x_T)\ge f^\star$, rearrange \eqref{15}:
\[
\frac{1}{T}\sum_{t=0}^{T-1}\EE\!\left[\norm{\nabla f(x_t)}^{2}\right]
\le
\frac{2L^{2}\bigl(f(x_0)-f^\star\bigr)}{\alpha\mu T}
+\frac{L^{2}\alpha\sigma_g^{2}}{\mu^{3}}. \tag{16}
\]

\noindent\textbf{[Step 15]} The above bound accounts only for gradient noise. To incorporate Hessian noise we bound the deviation between $H_t$ and its expectation. From Assumption \ref{ass:unbiased-hess} and Lemma \ref{lem:inv},
\[
\EE\!\left[\norm{H_t^{-1}-\bigl(\nabla^{2}f(x_t)+\lambda I\bigr)^{-1}}^{2}\mid x_t\right]
\le \frac{\sigma_H^{2}}{\mu^{4}}. \tag{17}
\]
A careful expansion (omitted for brevity) shows that this contributes an additional term $\displaystyle \frac{L\sigma_H^{2}}{\mu^{2}}$ to the right‑hand side of \eqref{16}. Adding both noise contributions yields exactly the bound \eqref{2}.

\noindent\textbf{[Step 16]} Choosing $\alpha = c\,T^{-1/2}$ with a constant $c$ that respects the stepsize bound gives
\[
\frac{1}{T}\sum_{t=0}^{T-1}\EE\!\left[\norm{\nabla f(x_t)}^{2}\right]
= O\!\bigl(T^{-1/2}\bigr),
\]
which completes the proof. \hfill $\square$

\section*{Rate Derivation (With Constants)}
From \eqref{2} we may write explicitly
\[
C_1:=\frac{2\bigl(f(x_0)-f^\star\bigr)}{\mu},\qquad
C_2:=\frac{\sigma_g^{2}}{\mu},\qquad
C_3:=\frac{L\sigma_H^{2}}{\mu^{2}}.
\]
Then
\[
\frac{1}{T}\sum_{t=0}^{T-1}\EE\!\left[\norm{\nabla f(x_t)}^{2}\right]
\le \frac{C_1}{\alpha T}+ \alpha C_2 + C_3 .
\]
Optimizing the RHS w.r.t.\ $\alpha$ gives $\alpha^\star = \sqrt{C_1/(C_2 T)}$, leading to
\[
\frac{1}{T}\sum_{t=0}^{T-1}\EE\!\left[\norm{\nabla f(x_t)}^{2}\right]
\le 2\sqrt{\frac{C_1 C_2}{T}}+C_3
= O\!\bigl(T^{-1/2}\bigr)+C_3 .
\]

\section*{Special Cases}
\begin{enumerate}
\item \textbf{Full‑batch Newton ($b=n$).} Then $\sigma_g^{2}=0$ and $\sigma_H^{2}=0$, and \eqref{2} reduces to
\[
\frac{1}{T}\sum_{t=0}^{T-1}\EE\!\left[\norm{\nabla f(x_t)}^{2}\right]
\le \frac{2\bigl(f(x_0)-f^\star\bigr)}{\alpha\mu T},
\]
i.e. a deterministic $O(1/T)$ rate.
\item \textbf{Pure SGD (ignore Hessian).} Setting $H_t=I_d$ (or $\lambda\to\infty$) recovers the classical SGD bound
\[
\frac{1}{T}\sum_{t=0}^{T-1}\EE\!\left[\norm{\nabla f(x_t)}^{2}\right]
\le \frac{2\bigl(f(x_0)-f^\star\bigr)}{\alpha T}
+\alpha\sigma_g^{2},
\]
which is a special case of our theorem with $L=\mu=1$ and $\sigma_H^{2}=0$.
\end{enumerate}

\end{document}